% GENERATED FILE. Edit canonical lessons, then run npm run generate:books.
% canonical-source-hash: fnv1a64:72ffe16e05a77058
% canonical-lessons: FR-C79-telephone, FR-C79-parapluie, FR-C79-argent, FR-C79-cadeau, FR-C79-lunettes

\chapter{Phone, Umbrella, Money, Present, Glasses}
\label{ch:fr-phone-umbrella-money-present-glasses}
\hlchaptermodality{79}

\begin{chapteropening}
\textbf{By the end of this chapter:} \emph{I can name a phone, an umbrella, money, a present, glasses.}

Complete the last lesson of chapter 79: I can name a phone, an umbrella, money, a present, glasses.
\end{chapteropening}

\section[le téléphone]{le téléphone — a phone}

\label{lesson:FR-C79-telephone}



\begin{quote}
Before the new one: say the French for a pencil, then the French for a bag.
\end{quote}

\subsection*{You'll want to know: le téléphone}
\textbf{le téléphone} — \textquotedblleft{}a phone\textquotedblright{}.

\begin{cousinweb}
\textbf{téléphone} — a 19th-century coinage from Greek \emph{tele}, \textquotedblleft{}far off\textquotedblright{}, and \emph{phone}, \textquotedblleft{}voice, sound\textquotedblright{}. (EtymDB 2.1)
\end{cousinweb}

\begin{grammarlens}[title={What you've built}]
One more word for this chapter.
\end{grammarlens}

\subsection*{Your turn}
\begin{itemize}
\raggedright
  \item \emph{Say it:} \emph{le téléphone}
  \item \emph{Say it:} \emph{le téléphone}, once more
  \item \emph{Say it:} say \emph{le téléphone}
  \item \emph{Recall:} say the French for a pencil, then the French for a bag, then say \emph{le téléphone} again
\end{itemize}

\begin{tcolorbox}[breakable,colback=teal!4,colframe=teal!35!black,title={Before you move on}]
What does le téléphone mean? (\textquotedblleft{}A phone\textquotedblright{}.) Say it once more.
\end{tcolorbox}
\section[le parapluie]{le parapluie — an umbrella}

\label{lesson:FR-C79-parapluie}



\begin{quote}
Before the new one: say the French for a bag, then the French for a phone.
\end{quote}

\subsection*{You'll want to know: le parapluie}
\textbf{le parapluie} — \textquotedblleft{}an umbrella\textquotedblright{}. Literally \textquotedblleft{}guard-rain\textquotedblright{}.

\begin{cousinweb}
\textbf{parapluie} — built from \emph{para-}, \textquotedblleft{}guarding against\textquotedblright{}, and the French word for \textquotedblleft{}rain\textquotedblright{}. \emph{para-} came from Italian \emph{parare}, \textquotedblleft{}to shield\textquotedblright{}, $\leftarrow$ Latin \emph{parare}, \textquotedblleft{}to prepare\textquotedblright{}. (EtymDB 2.1)
\end{cousinweb}

\begin{grammarlens}[title={What you've built}]
One more word for this chapter.
\end{grammarlens}

\subsection*{Your turn}
\begin{itemize}
\raggedright
  \item \emph{Say it:} \emph{le parapluie}
  \item \emph{Say it:} \emph{le parapluie}, once more
  \item \emph{Say it:} say \emph{le parapluie}
  \item \emph{Recall:} say the French for a bag, then the French for a phone, then say \emph{le parapluie} again
\end{itemize}

\begin{tcolorbox}[breakable,colback=teal!4,colframe=teal!35!black,title={Before you move on}]
What does le parapluie mean? (\textquotedblleft{}An umbrella\textquotedblright{}.) Say it once more.
\end{tcolorbox}
\section[l'argent]{l'argent — money; silver}

\label{lesson:FR-C79-argent}



\begin{quote}
Before the new one: say the French for a phone, then the French for an umbrella.
\end{quote}

\subsection*{You'll want to know: l'argent}
\textbf{l'argent} — \textquotedblleft{}money; silver\textquotedblright{}.

\begin{cousinweb}
\textbf{argent} $\leftarrow$ Latin \emph{argentum}, \textquotedblleft{}silver\textquotedblright{} — inherited. (EtymDB 2.1) Coins were silver, so the word came to mean money.
\end{cousinweb}

\begin{grammarlens}[title={What you've built}]
One more word for this chapter.
\end{grammarlens}

\subsection*{Your turn}
\begin{itemize}
\raggedright
  \item \emph{Say it:} \emph{l'argent}
  \item \emph{Say it:} \emph{l'argent}, once more
  \item \emph{Say it:} say \emph{l'argent}
  \item \emph{Recall:} say the French for a phone, then the French for an umbrella, then say \emph{l'argent} again
\end{itemize}

\begin{tcolorbox}[breakable,colback=teal!4,colframe=teal!35!black,title={Before you move on}]
What does l'argent mean? (\textquotedblleft{}Money; silver\textquotedblright{}.) Say it once more.
\end{tcolorbox}
\section[le cadeau]{le cadeau — a present, a gift}

\label{lesson:FR-C79-cadeau}



\begin{quote}
Before the new one: say the French for an umbrella, then the French for money; silver.
\end{quote}

\subsection*{You'll want to know: le cadeau}
\textbf{le cadeau} — \textquotedblleft{}a present, a gift\textquotedblright{}.

\begin{cousinweb}
\textbf{cadeau} $\leftarrow$ Old Provençal \emph{capdel}, \textquotedblleft{}a chief; a capital letter\textquotedblright{}, $\leftarrow$ Latin \emph{capitellum}, \textquotedblleft{}a little head\textquotedblright{} — borrowed. (EtymDB 2.1) It meant a flourished capital letter, then a flourish, a treat, and finally a gift.
\end{cousinweb}

\begin{grammarlens}[title={What you've built}]
One more word for this chapter.
\end{grammarlens}

\subsection*{Your turn}
\begin{itemize}
\raggedright
  \item \emph{Say it:} \emph{le cadeau}
  \item \emph{Say it:} \emph{le cadeau}, once more
  \item \emph{Say it:} say \emph{le cadeau}
  \item \emph{Recall:} say the French for an umbrella, then the French for money; silver, then say \emph{le cadeau} again
\end{itemize}

\begin{tcolorbox}[breakable,colback=teal!4,colframe=teal!35!black,title={Before you move on}]
What does le cadeau mean? (\textquotedblleft{}A present, a gift\textquotedblright{}.) Say it once more.
\end{tcolorbox}
\section[les lunettes]{les lunettes — glasses}

\label{lesson:FR-C79-lunettes}



\begin{quote}
Before the new one: say the French for money; silver, then the French for a present.
\end{quote}

\subsection*{You'll want to know: les lunettes}
\textbf{les lunettes} — \textquotedblleft{}glasses\textquotedblright{}.

\begin{cousinweb}
\textbf{lunettes} — built from \emph{lune}, \textquotedblleft{}moon\textquotedblright{}, with \emph{-ette}, \textquotedblleft{}little\textquotedblright{}: the round lenses were little moons. \emph{lune} $\leftarrow$ Latin \emph{luna}. (EtymDB 2.1)
\end{cousinweb}

\begin{grammarlens}[title={What you've built}]
One more word for this chapter.
\end{grammarlens}

\subsection*{Your turn}
\begin{itemize}
\raggedright
  \item \emph{Say it:} \emph{les lunettes}
  \item \emph{Say it:} \emph{les lunettes}, once more
  \item \emph{Say it:} say \emph{les lunettes}
  \item \emph{Recall:} say the French for money; silver, then the French for a present, then say \emph{les lunettes} again
\end{itemize}

\begin{tcolorbox}[breakable,colback=teal!4,colframe=teal!35!black,title={Before you move on}]
What does les lunettes mean? (\textquotedblleft{}Glasses\textquotedblright{}.) Say it once more.
\end{tcolorbox}
